\scanpage{39}\mspage{\deleted{\unsure{X bis|XI}}$\underline{\overline{\mathrm{XII}}}$}
On observe que $\Lambda^1$ et $\Lambda^2$ ont une arète en commun\inserted{, savoir,} $[s,s']$ et $\Lambda^1\cap\Lambda^2=[s,s']$.

\uline{Lemme} XIV. Si $D_s$ est un disque satisfaisant b) de la definition 2 par rapport à $\Lambda^1$ et à $\Lambda^2$ on a que $\operatorname{int}\Lambda^1\supset\Lambda^2\cap D_s-[s,s']\cap D_s$ ou bien $\operatorname{int}\Lambda^2\supset\Lambda^1\cap D_s\deleted{\unsure{-|\subset}}[s,s']\cap D_s$ \deleted{\unsure{(}}et resp. \inserted{que} $\operatorname{ext}\Lambda^2\supset\Lambda^1\cap D_s-[s,s']\cap D_s$ ou bien $\operatorname{ext}\Lambda^1\supset\Lambda^2\cap D_s-[s,s']\cap D_s$\deleted{)}.

\uline{Demonstration} : D'aprés le lemme XIII, $\operatorname{ext}\Lambda^1\cap D_s\supset\operatorname{ext}\partial E(\Gamma)\cap D_s$, $\operatorname{ext}\Lambda^2\cap D_s\supset\operatorname{ext}\partial E(\Gamma)\cap D_s$ donc $\operatorname{ext}\Lambda^1\cap\operatorname{ext}\Lambda^2\cap D_s\ne\varnothing$.

Comme $\Lambda^1\cap D_s$ et $\Lambda^2\cap D_s$ ont un coté en commun il résulte que le \inserted{secteur angulaire} $\operatorname{ext}\Lambda^1\cap D_s$ contient ou est contenu dans l'un des secteurs angulaires determinés par $\Lambda^2\cap D_s$. Mais $\operatorname{ext}\Lambda^1\cap\operatorname{ext}\Lambda^2\cap D_s\ne\varnothing$ entraîne que ce secteur est précisement $\operatorname{ext}\Lambda^2\cap D_s$, cela implique (sans peine)
\[
\operatorname{int}\Lambda^1\supset\Lambda^2\cap D_s-[s,s']\cap D_s
\quad\text{ou bien}\quad
\deleted{\operatorname{ext}}\operatorname{int}\Lambda^2\supset\Lambda^1\cap D_s-[s,s']\cap D_s,
\]
c.q.f.d.

La conjonction du lemme $\underline{\overline{\mathrm{XII}}}$ et du lemme $\underline{\overline{\mathrm{XIII}}}$ \deleted{\unsure{entraîne|entraînent}} montre que
\[
\begin{aligned}
\operatorname{int}\Lambda^1\supset{}&\Lambda^2-\{g,s'\}
&&\text{ou bien }\operatorname{int}\Lambda^2\supset\Lambda^1-\{s,s'\}\ \text{et on a}\\[-2pt]
&\inserted{\substack{\text{''}\\\Gamma^2_{s,s'}-\{s,s'\}}}
&&\inserted{-\qquad\substack{\text{''}\\\Gamma^{\unsure{1|2}}_{s,s'}.-\{s,s'\}}}\\
\text{respectivement}\quad&\operatorname{ext}\Lambda^2\supset\Lambda^1-\{s,s'\}
&&\text{et }\operatorname{ext}\Lambda^1\supset\Lambda^2-\{s,s'\},
\end{aligned}
\]

Supposons pour fixer les idées que $\operatorname{int}\Lambda^1\supset\Lambda^2-\{s,s'\}$ donc aussi que $\operatorname{ext}\Lambda^2\supset\Lambda^1-\{s,s'\}$.

On a donc :

\uline{Lemme XV} : $\operatorname{int}\Lambda^2\subset\operatorname{int}\Lambda^1$.
% archive sheet number: 38
